\documentclass[10pt,a4paper]{amsart}
\usepackage{amsmath,amssymb,amsthm,graphicx}
% Retained minimal preamble for the Stacks Project chapter sites-cohomology.tex (Cohomology on Sites).
% Verbatim excerpt of the chapter preamble: Xy-pic setup (native preamble lines 14-18) and the native
% theorem declarations with the mathematical macros (native preamble lines 154-217). Copied mechanically;
% nothing added, removed or redefined. The wrapper class and the amsmath/amssymb/amsthm/graphicx base
% packages are supplied by the assembler, so the class/externaldocument part of the native preamble is omitted.
% For commutative diagrams we use Xy-pic
\usepackage[all]{xy}

% We use 2cell for 2-commutative diagrams.
\xyoption{2cell}
% Theorem environments.
%
\theoremstyle{plain}
\newtheorem{theorem}[subsection]{Theorem}
\newtheorem{proposition}[subsection]{Proposition}
\newtheorem{lemma}[subsection]{Lemma}

\theoremstyle{definition}
\newtheorem{definition}[subsection]{Definition}
\newtheorem{example}[subsection]{Example}
\newtheorem{exercise}[subsection]{Exercise}
\newtheorem{situation}[subsection]{Situation}

\theoremstyle{remark}
\newtheorem{remark}[subsection]{Remark}
\newtheorem{remarks}[subsection]{Remarks}

\numberwithin{equation}{subsection}

% Macros
%
\def\lim{\mathop{\mathrm{lim}}\nolimits}
\def\colim{\mathop{\mathrm{colim}}\nolimits}
\def\Spec{\mathop{\mathrm{Spec}}}
\def\Hom{\mathop{\mathrm{Hom}}\nolimits}
\def\Ext{\mathop{\mathrm{Ext}}\nolimits}
\def\SheafHom{\mathop{\mathcal{H}\!\mathit{om}}\nolimits}
\def\SheafExt{\mathop{\mathcal{E}\!\mathit{xt}}\nolimits}
\def\Sch{\mathit{Sch}}
\def\Mor{\mathop{\mathrm{Mor}}\nolimits}
\def\Ob{\mathop{\mathrm{Ob}}\nolimits}
\def\Sh{\mathop{\mathit{Sh}}\nolimits}
\def\NL{\mathop{N\!L}\nolimits}
\def\CH{\mathop{\mathrm{CH}}\nolimits}
\def\proetale{{pro\text{-}\acute{e}tale}}
\def\etale{{\acute{e}tale}}
\def\QCoh{\mathit{QCoh}}
\def\Ker{\mathop{\mathrm{Ker}}}
\def\Im{\mathop{\mathrm{Im}}}
\def\Coker{\mathop{\mathrm{Coker}}}
\def\Coim{\mathop{\mathrm{Coim}}}

% Boxtimes
%
\DeclareMathSymbol{\boxtimes}{\mathbin}{AMSa}{"02}

%
% Macros for moduli stacks/spaces
%
\def\QCohstack{\mathcal{QC}\!\mathit{oh}}
\def\Cohstack{\mathcal{C}\!\mathit{oh}}
\def\Spacesstack{\mathcal{S}\!\mathit{paces}}
\def\Quotfunctor{\mathrm{Quot}}
\def\Hilbfunctor{\mathrm{Hilb}}
\def\Curvesstack{\mathcal{C}\!\mathit{urves}}
\def\Polarizedstack{\mathcal{P}\!\mathit{olarized}}
\def\Complexesstack{\mathcal{C}\!\mathit{omplexes}}
% \Pic is the operator that assigns to X its picard group, usage \Pic(X)
% \Picardstack_{X/B} denotes the Picard stack of X over B
% \Picardfunctor_{X/B} denotes the Picard functor of X over B
\def\Pic{\mathop{\mathrm{Pic}}\nolimits}
\def\Picardstack{\mathcal{P}\!\mathit{ic}}
\def\Picardfunctor{\mathrm{Pic}}
\def\Deformationcategory{\mathcal{D}\!\mathit{ef}}
\usepackage{cleveref}
\crefname{originalreference}{original reference}{original references}
\Crefname{originalreference}{Original reference}{Original references}
\makeatletter
\newlabel{lemma-local-actual}{{08C9}{source}{}{sourcepointer-1}{}}
\newlabel{lemma-local-actual@cref}{{[originalreference][0][]08C9}{[1][0][]source}{}{}{}}
\makeatother
\crefname{originaldefinition}{original definition}{original definitions}
\Crefname{originaldefinition}{Original definition}{Original definitions}
\makeatletter
\newlabel{sites-definition-site}{{00VH}{source}{}{sourcepointer-2}{}}
\newlabel{sites-definition-site@cref}{{[originaldefinition][0][]00VH}{[1][0][]source}{}{}{}}
\makeatother
\title{Pseudo-coherence of an object is independent of the representing complex and can be checked locally}
\author{The Stacks Project Authors}
\date{stacks-project:sites-cohomology.tex:local-snapshot-20261011; GNU FDL 1.2 or later}
\begin{document}
\maketitle
\noindent\textit{Editorial note.} All delivered mathematical text is the author's own, copied from the pinned Stacks Project chapter file sites-cohomology.tex; no statement or proof step was rewritten, reordered or removed.\par\medskip
\noindent\textit{Editorial note.} The counted proof is the author's complete argument: the proof is written out in full: the refinement and representation of the local maps by complexes, the deduction of pseudo-coherence for an arbitrary representative, the passage from localisations to the global statement, and the local-to-global step for part (3).\par\medskip
\noindent\textit{Editorial note.} The label lemma-pseudo-coherent-independent-representative names the counted statement, and the proof environment immediately following it in the source is delivered verbatim as its proof.\par\medskip
\noindent\textit{Editorial note.} The counted statement is carried by the author's own lemma environment in the chapter source.\par\medskip
\noindent\textit{Editorial note.} The statement and proof are copied from the pinned chapter file; the two cited inputs are kept as external pointers with their canonical Stacks tags.\par\medskip
\noindent\textit{Editorial note.} The preamble is the chapter's own retained minimal preamble (sources/stacks-sites-cohomology/retained-preamble.tex), a verbatim excerpt of the native preamble: the Xy-pic setup and the native theorem declarations with the mathematical macros.\par\medskip
\noindent\textit{Allowed standard prerequisites (editorial).}\begin{itemize}
\item m-pseudo-coherent complexes of modules on a ringed site. Same chapter, the definition of m-pseudo-coherence through a strictly perfect complex mapping to the object with cohomology isomorphisms above m and a surjection in degree m. Supplies the property verified for the chosen representative and for its localisations.
\item Local representability of maps in the derived category. Same chapter, the lemma allowing a map in the derived category to be represented by a map of complexes after refining a covering. Supplies the passage from the maps alpha\_i in the derived category to maps of complexes on a refinement.
\item The definition of a site through its coverings. Sites, the definition of a site, in particular the part listing the axioms of coverings. Supplies the local-to-global step in part (3).
\end{itemize}\medskip
\ifdefined\hypertarget\hypertarget{sourcepointer-1}{}\fi
\noindent\textit{Source reference.} Original reference 08C9: sources/stacks-sites-cohomology/pinned-original-source.tex, source lines 11881--11881. Prerequisite pointer invoked at proof line(s) 12089; it names a standard fact of the same chapter that the counted argument uses, and every step of the counted proof built on it is retained verbatim in the delivered body.\par\medskip
\ifdefined\hypertarget\hypertarget{sourcepointer-2}{}\fi
\noindent\textit{Source reference.} Original definition 00VH: sources/stacks-chow/tags.txt, source lines 1147--1147. Prerequisite pointer invoked at proof line(s) 12106; it names a standard fact of a sibling chapter that the counted argument uses, and every step of the counted proof built on it is retained verbatim in the delivered body.\par\medskip
\begin{lemma}
\label{lemma-pseudo-coherent-independent-representative}
Let $(\mathcal{C}, \mathcal{O})$ be a ringed site.
Let $E$ be an object of $D(\mathcal{O})$.
\begin{enumerate}
\item If $\mathcal{C}$ has a final object $X$ and if there exist a covering
$\{U_i \to X\}$, strictly perfect complexes $\mathcal{E}_i^\bullet$ of
$\mathcal{O}_{U_i}$-modules, and
maps $\alpha_i : \mathcal{E}_i^\bullet \to E|_{U_i}$ in
$D(\mathcal{O}_{U_i})$ with $H^j(\alpha_i)$ an isomorphism for $j > m$
and $H^m(\alpha_i)$ surjective, then $E$ is $m$-pseudo-coherent.
\item If $E$ is $m$-pseudo-coherent, then any complex of $\mathcal{O}$-modules
representing $E$ is $m$-pseudo-coherent.
\item If for every object $U$ of $\mathcal{C}$ there exists a covering
$\{U_i \to U\}$ such that $E|_{U_i}$ is $m$-pseudo-coherent, then
$E$ is $m$-pseudo-coherent.
\end{enumerate}
\end{lemma}

\begin{proof}
Let $\mathcal{F}^\bullet$ be any complex representing $E$
and let $X$, $\{U_i \to X\}$, and $\alpha_i : \mathcal{E}_i \to E|_{U_i}$
be as in (1). We will show that $\mathcal{F}^\bullet$ is $m$-pseudo-coherent
as a complex, which will prove (1) and (2) in case $\mathcal{C}$ has a
final object. By Lemma \ref{lemma-local-actual}
we can after refining the covering $\{U_i \to X\}$
represent the maps $\alpha_i$ by maps of complexes
$\alpha_i : \mathcal{E}_i^\bullet \to \mathcal{F}^\bullet|_{U_i}$.
By assumption
$H^j(\alpha_i)$ are isomorphisms for $j > m$, and $H^m(\alpha_i)$
is surjective whence $\mathcal{F}^\bullet$ is
$m$-pseudo-coherent.

\medskip\noindent
Proof of (2). By the above we see that $\mathcal{F}^\bullet|_U$ is
$m$-pseudo-coherent as a complex of $\mathcal{O}_U$-modules for all
objects $U$ of $\mathcal{C}$. It is a formal consequence of the definitions
that $\mathcal{F}^\bullet$ is $m$-pseudo-coherent.

\medskip\noindent
Proof of (3). Follows from the definitions and
Sites, Definition \ref{sites-definition-site} part (2).
\end{proof}

\end{document}
